\section{Discussion}\label{sec:discussion}

\subsection{Practical relevance}
The bound of Theorem~\ref{thm:main} is far from a practical running
time. In \eqref{eq:explicitwork}, the net size $J_*$ and the grid factor
$(K+2)^{w+1}$ are $(w+1)$st powers of quantities that grow polynomially
with $n$, $C$, $H$, $\mu^{-1}$, $R$, and $\sigma^{-1}$. To see the
scale, take $n=100$, $w=2$, $H=2$, $\mu=C=R=1$, and $\sigma=1/10$. Then
\eqref{eq:mainparameters} and \eqref{eq:gridB} give $J_*>5\cdot10^{20}$
and $K>6\cdot10^4$, so $J_*(K+2)^{w+1}>10^{35}$ even before the factors
$b$ and $P_w(\ell)$ are included. Moreover, every oracle call runs a
full grid dynamic program, and each message uses at most $J_*(1+en)$
oracle calls in expectation (Step~4 of the proof of
Theorem~\ref{thm:main}). These factors come from bounds chosen for the
analysis: one Lipschitz constant and one grid serve all supports, and the
net ignores where the optimal support changes. We have not tried to
reduce them. An exact-rational reference implementation, available from
the authors, replaces $X$ and $L$ by bounds specific to each message. We
used it only to check correctness on small instances, including supports
that are optimal only at a tie point.

Nor does Theorem~\ref{thm:main} imply a speedup over existing exact
methods, such as the parametric dynamic programs of
\cite{BhathenaTrees,BhathenaStructured} discussed in
Section~\ref{sec:related}. In the synthetic banded experiments of
\cite[Section~7.2, Table~1]{BhathenaStructured}, with bandwidth two or
four and $n$ up to $2{,}000$, the reported average number of quadratic
pieces retained after pruning (branches, in our terminology) is between
23 and $1{,}139$. Exponential messages such as that of
Theorem~\ref{thm:fixed-data} therefore need not arise on such data. The
guarantee of Theorem~\ref{thm:main} matters most for instances with many
near ties, like the example of that theorem.

The grid oracle of Lemma~\ref{lem:gridoracle} solves the problem
without side constraints. A cardinality bound on the support can be
added at a polynomial cost, by also recording the number of active
vertices in the dynamic-program tables. Nonlinear constraints on $x$
would need a different oracle. Theorem~\ref{thm:dictionary} itself
applies to any support family for which a restricted additive oracle is
available.

\subsection{Possible practical uses}
Two parts of the analysis may still be useful in practice, although we
have not tested either of them.

The first is the penalty perturbation itself, used inside a pruned exact
dynamic program. The exact algorithm of \cite{BhathenaStructured} prunes
branches that are never optimal on a bounded region of boundary values.
Its complexity bound rests on a deterministic margin condition
\cite[Definition~5 and Theorem~1]{BhathenaStructured}. At each bag
$B_u$, give every support of the local subproblem the optimal cost of
that subproblem with the bag variables set to zero, and group the
supports into classes by their indicator values on the vertices within a
fixed graph distance of $B_u$. The $\eta$-optimal set of a class
consists of its supports whose cost is within $\eta$ of the least cost
in the class. The condition requires a single $k$ that bounds the size
of every $\eta$-optimal set, over all bags and classes. Under our
perturbation, Lemma~\ref{lem:count} bounds the expected size of each
such set by $e$ when $\eta\le\sigma/(2n)$. An expected bound for each
set is weaker than their uniform bound, so this gives no running-time
guarantee for their algorithm. Perturbing the penalties may still
shorten the lists that their algorithm keeps on instances with many near
ties. The loss in accuracy is bounded in advance: from an exact solution
of the perturbed problem, \eqref{eq:certificate} gives lower and upper
bounds on the original optimum that differ by at most $n\sigma$ for
every draw.

The second is certified enumeration with oracles other than the grid
dynamic program. Lemma~\ref{lem:enumeration} works with any restricted
additive oracle in the sense of \eqref{eq:oracle}. From such an oracle
it builds a list $E$ that contains every support within $\delta$ of the
optimum and only supports within $\delta+2\varepsilon$. It makes at most
$1+n|E|$ oracle calls, so its cost grows with the number of near-optimal
supports, which can be exponential without perturbation.

Approximate decision diagrams, described in Section~\ref{sec:related},
can serve as such oracles on the structures they cover. In the truncated
diagrams of \cite{GomezBanded} and the neighborhood-merged diagrams of
\cite{ChoiDiagrams}, each support is encoded by a path whose length
differs from the cost of that support by at most a bound that is the
same for all supports. For \cite{GomezBanded} the bound follows by
summing \cite[Proposition~7]{GomezBanded} along the path; for
\cite{ChoiDiagrams} it is \cite[Proposition~9]{ChoiDiagrams}, and it is
proportional to the merging tolerance. If the truncation depth or the
merging tolerance is chosen so that this bound is at most
$\varepsilon/2$, a shortest path through the arcs consistent with the
fixed bits, returned with the exact cost of its support, is a restricted
additive oracle with accuracy $\varepsilon$. The diagrams also need not
be rebuilt at each net point. As noted in Section~\ref{sec:related},
they are built from the matrix alone, and only their arc lengths change
with the linear coefficients and penalties. In \eqref{eq:message} the
boundary value $t$ enters only through the linear coefficients
$c_I+2Q_{IS}t$. Hence one diagram for $Q_{II}$, with its accuracy set for
the largest linear coefficients on the box, serves every net point of a
message.

\subsection{Open questions}
\begin{enumerate}
\item \emph{Computational study.} How large are the dictionaries of
perturbed instances in practice, and how do they compare with the pruned
lists of exact parametric dynamic programs
\cite{BhathenaTrees,BhathenaStructured} on the same data?
\item \emph{Avoiding the parameter net.} The net covers the whole box
uniformly, and its size $J_*$ is a large factor in
\eqref{eq:explicitwork}. Is there an adaptive search over boundary values
with a provable expected bound, for example one that refines the net
only where the enumerated lists change?
\item \emph{Recursive composition.} Our algorithm computes each message
directly from its own subproblem. This avoids products of child
dictionaries but repeats work. Can child dictionaries be combined
bottom-up, with an expected bound on the intermediate lists, without
additional assumptions on $Q$?
\item \emph{Deterministic conditions.} Theorem~\ref{thm:main} bounds the
expected dictionary size over perturbation draws for every instance, but
it does not say which unperturbed instances have small messages. Which
checkable conditions on unperturbed data, other than the margin and
volume-growth conditions of \cite{BhathenaStructured}, guarantee
polynomial message size at fixed treewidth?
\item \emph{Finer representation lower bounds.} The star argument of
Appendix~\ref{sec:representation} concerns the inverse-principal polytope
including its inverse entries. Are there exponential lower bounds for
formulations that project out the inverse entries, such as extended
formulations of the closed convex hull of the epigraph?
\end{enumerate}
